% Appendix A: Key Lean 4 Declarations
%
% This appendix presents the formal Lean 4 definitions and proofs of the
% kernel's key contributions. Each code block is extracted verbatim from
% the kernel source. The reader can verify any declaration by cloning the
% repository and running `lake build`.

\chapter{Key Lean 4 Declarations}\label{app:lean-code}

This appendix displays the formal Lean 4 code for the definitions and
theorems claimed in the preceding chapters. Every declaration compiles
against mathlib4 \texttt{fde0cc5} with zero \texttt{sorry}.

\section{Core types}

\begin{lstlisting}[caption={Concept and ConceptClass (Basic.lean)}]
def Concept (X : Type u) (Y : Type v) := X -> Y

abbrev ConceptClass (X : Type u) (Y : Type v) := Set (Concept X Y)
\end{lstlisting}

\begin{lstlisting}[caption={BatchLearner (Learner/Core.lean)}]
structure BatchLearner (X : Type u) (Y : Type v) where
  hypotheses : HypothesisSpace X Y
  learn : {m : N} -> (Fin m -> X * Y) -> Concept X Y
  output_in_H : forall {m : N} (S : Fin m -> X * Y), learn S (*@$\in$@*) hypotheses
\end{lstlisting}

\section{Combinatorial dimensions}

\begin{lstlisting}[caption={Shattering and VC dimension (Complexity/VCDimension.lean)}]
def Shatters (X : Type u) (C : ConceptClass X Bool) (S : Finset X) : Prop :=
  forall f : S -> Bool, exists c (*@$\in$@*) C, forall x : S, c (x : X) = f x

noncomputable def VCDim (X : Type u) (C : ConceptClass X Bool) : WithTop N :=
  iSup (S : Finset X) (_ : Shatters X C S), (S.card : WithTop N)
\end{lstlisting}

\begin{lstlisting}[caption={Littlestone dimension (Complexity/GameInfra.lean)}]
noncomputable def LittlestoneDim (X : Type) (C : ConceptClass X Bool) :
    WithBot (WithTop N) :=
  iSup (n : N) (_ : exists T : LTree X n, T.isShattered C),
    ((n : WithTop N) : WithBot (WithTop N))
\end{lstlisting}

\section{Learnability predicates}

\begin{lstlisting}[caption={PAC learnability (Criterion/PAC.lean)}]
def PACLearnable (X : Type u) [MeasurableSpace X]
    (C : ConceptClass X Bool) : Prop :=
  exists (L : BatchLearner X Bool) (mf : R -> R -> N),
    forall (e d : R), 0 < e -> 0 < d ->
      forall (D : Measure X), IsProbabilityMeasure D ->
        forall (c : Concept X Bool), c (*@$\in$@*) C ->
          let m := mf e d
          Measure.pi (fun _ : Fin m => D)
            { xs : Fin m -> X |
              D { x | L.learn (fun i => (xs i, c (xs i))) x (*@$\neq$@*) c x }
                (*@$\leq$@*) ENNReal.ofReal e }
            (*@$\geq$@*) ENNReal.ofReal (1 - d)
\end{lstlisting}

\begin{lstlisting}[caption={Online learnability (Criterion/Online.lean)}]
def OnlineLearnable (X : Type u) (Y : Type v) [DecidableEq Y]
    (C : ConceptClass X Y) : Prop :=
  exists (M : N), MistakeBounded X Y C M
\end{lstlisting}

\begin{lstlisting}[caption={EX-learnability, Gold-style (Criterion/Gold.lean)}]
def EXLearnable (X : Type u) (C : ConceptClass X Bool) : Prop :=
  exists (L : GoldLearner X Bool),
    forall (c : Concept X Bool), c (*@$\in$@*) C ->
      forall (T : TextPresentation X c),
        exists (t0 : N), forall (t : N), t (*@$\geq$@*) t0 ->
          L.conjecture (dataUpTo T.toDataStream t) = c
\end{lstlisting}

\section{The fundamental theorem (5-way equivalence)}

\begin{lstlisting}[caption={Fundamental theorem of statistical learning (Theorem/PAC.lean:293)}]
theorem fundamental_theorem (X : Type u) [MeasurableSpace X]
    [MeasurableSingletonClass X]
    (C : ConceptClass X Bool)
    [MeasurableConceptClass X C] :
    (PACLearnable X C (*@$\leftrightarrow$@*) VCDim X C < (*@$\top$@*)) (*@$\wedge$@*)
    ((VCDim X C < (*@$\top$@*)) (*@$\leftrightarrow$@*)
      exists (k : N) (cs : CompressionSchemeWithInfo0 X Bool C), cs.size = k) (*@$\wedge$@*)
    ((VCDim X C < (*@$\top$@*)) (*@$\leftrightarrow$@*)
      forall e > 0, exists m0, forall (D : Measure X),
        IsProbabilityMeasure D ->
        forall m (*@$\geq$@*) m0, RademacherComplexity X C D m < e) (*@$\wedge$@*)
    (PACLearnable X C ->
      exists (L : BatchLearner X Bool) (mf : R -> R -> N),
        -- sample complexity sandwich + lower bound
        ...) (*@$\wedge$@*)
    ((VCDim X C < (*@$\top$@*)) (*@$\leftrightarrow$@*)
      exists (d : N), forall (m : N), d (*@$\leq$@*) m ->
        GrowthFunction X C m (*@$\leq$@*)
          Finset.sum (Finset.range (d + 1)) (Nat.choose m)) :=
  -- Proof: assembles from component theorems
  (*@$\langle$@*)vc_characterization X C,
   fundamental_vc_compression X C,
   (vc_characterization X C).symm.trans (fundamental_rademacher X C),
   pac_sample_complexity_sandwich X C,
   (*@$\langle$@*)vcdim_finite_imp_growth_bounded X C,
    growth_bounded_imp_vcdim_finite X C(*@$\rangle\rangle$@*)
\end{lstlisting}

\section{Characterization theorems}

\begin{lstlisting}[caption={Littlestone characterization (Theorem/Online.lean:419)}]
theorem littlestone_characterization (X : Type)
    (C : ConceptClass X Bool) :
    OnlineLearnable X Bool C (*@$\leftrightarrow$@*) LittlestoneDim X C < (*@$\top$@*) :=
  (*@$\langle$@*)forward_direction X C, backward_direction X C(*@$\rangle$@*)
\end{lstlisting}

\begin{lstlisting}[caption={Gold's theorem (Theorem/Gold.lean:19, first 30 lines of 200)}]
theorem gold_theorem (X : Type u) [Countable X] [DecidableEq X]
    (C : ConceptClass X Bool)
    (h_finite : forall (S : Finset X), (fun x => decide (x (*@$\in$@*) S)) (*@$\in$@*) C)
    (h_infinite : exists c (*@$\in$@*) C, Set.Infinite { x | c x = true }) :
    (*@$\neg$@*) EXLearnable X C := by
  intro (*@$\langle$@*)L, hL(*@$\rangle$@*)
  obtain (*@$\langle$@*)c_inf, hc_inf_mem, hc_inf_inf(*@$\rangle$@*) := h_infinite
  let phi := hc_inf_inf.natEmbedding
  have hpos_ne : Nonempty { x : X | c_inf x = true } :=
    hc_inf_inf.nonempty.to_subtype
  obtain (*@$\langle$@*)enum, henum(*@$\rangle$@*) :=
    exists_surjective_nat ({ x : X | c_inf x = true })
  let S : N -> Finset X := fun n =>
    ((Finset.range (n + 1)).image (fun i => (phi i).val)) (*@$\cup$@*)
    ((Finset.range (n + 1)).image (fun i => (enum i).val))
  -- ... (152 more lines: constructs locking sequence via
  --      interleaving enumeration of c_inf's support with
  --      finite-set indicators, derives contradiction)
\end{lstlisting}

\section{Paradigm separations}

\begin{lstlisting}[caption={PAC does not imply Online (Theorem/Separation.lean:1175)}]
theorem pac_not_implies_online :
    exists (X : Type) (_ : MeasurableSpace X) (C : ConceptClass X Bool),
      PACLearnable X C (*@$\wedge$@*) (*@$\neg$@*) OnlineLearnable X Bool C := by
  refine (*@$\langle$@*)N, (*@$\top$@*), thresholdClass, ?_, ?_(*@$\rangle$@*)
  . -- PAC learnable: VCDim = 1 < (*@$\top$@*), via UC route
    exact vcdim_finite_imp_pac_via_uc' N thresholdClass
      vcdim_threshold_finite ...
  . -- (*@$\neg$@*) OnlineLearnable: LittlestoneDim = (*@$\top$@*)
    intro hol
    have hfin := forward_direction N thresholdClass hol
    rw [ldim_threshold_top] at hfin
    exact lt_irrefl _ hfin
\end{lstlisting}

\begin{lstlisting}[caption={Online implies PAC (Theorem/Separation.lean:132)}]
theorem online_imp_pac (X : Type u) [MeasurableSpace X]
    (C : ConceptClass X Bool) (hol : OnlineLearnable X Bool C)
    [MeasurableConceptClass X C] :
    PACLearnable X C := by
  obtain (*@$\langle$@*)M, hM(*@$\rangle$@*) := hol
  have hvcdim : VCDim X C < (*@$\top$@*) := by
    calc VCDim X C (*@$\leq$@*) M := vcdim_le_of_mistake_bounded hM
      _ < (*@$\top$@*) := WithTop.coe_lt_top M
  exact vcdim_finite_imp_pac_via_uc' X C hvcdim ...
\end{lstlisting}

\section{Measurability layer (kernel's novel contribution)}

\begin{lstlisting}[caption={WellBehavedVCMeasTarget (Complexity/Measurability.lean:388)}]
def WellBehavedVCMeasTarget
    (X : Type u) [MeasurableSpace X]
    (C : ConceptClass X Bool) : Prop :=
  forall (D : Measure X) [IsProbabilityMeasure D]
    (c : Concept X Bool), Measurable c ->
    forall (m : N) (e : R),
      NullMeasurableSet
        {p : (Fin m -> X) * (Fin m -> X) | exists h (*@$\in$@*) C,
          EmpiricalError X Bool h
            (fun i => (p.2 i, c (p.2 i))) (zeroOneLoss Bool) -
          EmpiricalError X Bool h
            (fun i => (p.1 i, c (p.1 i))) (zeroOneLoss Bool) (*@$\geq$@*) e / 2}
        ((Measure.pi (fun _ : Fin m => D)).prod
         (Measure.pi (fun _ : Fin m => D)))
\end{lstlisting}

\begin{lstlisting}[caption={KrappWirthWellBehaved (Complexity/Measurability.lean:241)}]
class KrappWirthWellBehaved (X : Type u) [MeasurableSpace X]
    (C : ConceptClass X Bool) : Prop extends MeasurableHypotheses X C where
  V_measurable : KrappWirthV X C
  U_measurable : KrappWirthU X C
\end{lstlisting}

\begin{lstlisting}[caption={Borel-analytic separation witness (Theorem/BorelAnalyticSeparation.lean:207)}]
theorem singleton_badEvent_not_measurable
    (A : Set R) (hA_non : (*@$\neg$@*) MeasurableSet A) :
    (*@$\neg$@*) MeasurableSet (singletonBadEvent A) := by
  intro hbad
  rw [singleton_badEvent_eq_preimage_planar A] at hbad
  have hmeas : Measurable samplePair1ToPlane :=
    Measurable.prod
      ((measurable_pi_apply 0).comp measurable_fst)
      ((measurable_pi_apply 0).comp measurable_snd)
  have hsurj : Function.Surjective samplePair1ToPlane := by
    intro (*@$\langle$@*)x, y(*@$\rangle$@*)
    exact (*@$\langle$@*)(fun _ => x, fun _ => y), by simp [samplePair1ToPlane](*@$\rangle$@*)
  have hplanar :=
    (hmeas.measurableSet_preimage_iff_of_surjective hsurj).mp hbad
  exact planarWitnessEvent_not_measurable A hA_non hplanar
\end{lstlisting}

\begin{lstlisting}[caption={Analytic sets are null-measurable (PureMath/AnalyticMeasurability.lean:91)}]
theorem analyticSet_nullMeasurableSet
    {(*@$\alpha$@*) : Type*}
    [TopologicalSpace (*@$\alpha$@*)] [MeasurableSpace (*@$\alpha$@*)]
    [BorelSpace (*@$\alpha$@*)] [PolishSpace (*@$\alpha$@*)]
    {(*@$\mu$@*) : Measure (*@$\alpha$@*)} [IsFiniteMeasure (*@$\mu$@*)]
    {s : Set (*@$\alpha$@*)} (hs : AnalyticSet s) :
    NullMeasurableSet s (*@$\mu$@*) := by
  classical
  obtain (*@$\langle$@*)t, hst, ht_meas, ht_eq(*@$\rangle$@*) :=
    exists_measurable_superset (*@$\mu$@*) s
  have hzero : (*@$\mu$@*) (t \ s) = 0 := by
    by_contra hne
    have hfin_diff : (*@$\mu$@*) (t \ s) (*@$\neq$@*) (*@$\top$@*) := measure_ne_top (*@$\mu$@*) _
    have hpos : 0 < (*@$\mu$@*).real (t \ s) :=
      ENNReal.toReal_pos hne hfin_diff
    obtain (*@$\langle$@*)K, hKc, hKs, hKapprox(*@$\rangle$@*) :=
      hs.exists_isCompact_measureReal_gt
        ((*@$\mu$@*) := (*@$\mu$@*)) ((*@$\mu$@*).real (t \ s) / 2) (by positivity)
    have hKt : K (*@$\subseteq$@*) t := fun x hx => hst (hKs hx)
    have hKmeas : MeasurableSet K := hKc.isClosed.measurableSet
    have hdiff_eq : (*@$\mu$@*).real (t \ K) = (*@$\mu$@*).real t - (*@$\mu$@*).real K :=
      measureReal_diff hKt hKmeas
    have ht_real : (*@$\mu$@*).real t = (*@$\mu$@*).real s := by
      simp only [Measure.real]; rw [ht_eq]
    have hsub : t \ s (*@$\subseteq$@*) t \ K :=
      Set.diff_subset_diff_right hKs
    have hle : (*@$\mu$@*).real (t \ s) (*@$\leq$@*) (*@$\mu$@*).real (t \ K) :=
      measureReal_mono hsub
    linarith
  have h_ae : s =(*@${}^{\text{ae}}$@*)[(*@$\mu$@*)] t := by
    rw [Filter.eventuallyEq_comm, ae_eq_set]
    exact (*@$\langle$@*)hzero, by simp [Set.diff_eq_empty.mpr hst](*@$\rangle$@*)
  exact ht_meas.nullMeasurableSet.congr h_ae.symm
\end{lstlisting}
